% Generated draft: requires mathematical and source review.
\documentclass[11pt]{article}
\usepackage[margin=25mm]{geometry}
\usepackage{fontspec}
\setmainfont{Latin Modern Roman}
\usepackage{amsmath,amssymb,mathtools,unicode-math}
\setmathfont{Latin Modern Math}
\usepackage{enumitem,needspace,xurl,hyperref}
\setlength{\parindent}{0pt}
\setlength{\parskip}{5pt}
\setlength{\emergencystretch}{3em}
\newcommand{\N}{\mathbb{N}}
\newcommand{\Z}{\mathbb{Z}}
\newcommand{\Q}{\mathbb{Q}}
\newcommand{\R}{\mathbb{R}}
\newcommand{\C}{\mathbb{C}}
\newcommand{\F}{\mathbb{F}}
\newcommand{\A}{\mathbb{A}}
\newcommand{\PP}{\mathbb{P}}
\newcommand{\E}{\mathbb{E}}
\newcommand{\eps}{\varepsilon}
\newcommand{\dd}{\,\mathrm{d}}
\begin{document}
\section*{UnsolvedMath ID: 30000510}
\textbf{Source problem code:} catalogue entry
\subsection{Definitions and mathematical statement}
For an integer \(n\geq3\), give
\[
U_n=\{(x_0,\ldots,x_{n-1})\in(\R\mathrm P^1)^n:
 x_i\ne x_{i+1}\text{ for all }i\pmod n,\ |\{x_0,\ldots,x_{n-1}\}|\geq3\}
\]
the subspace topology from the product of projective lines. The final adjacency condition includes \(x_{n-1}\ne x_0\). Nonadjacent coordinates may coincide. The indices are labeled: no cyclic rotation, reflection, or relabeling quotient is imposed.

Let \(\operatorname{SL}_2(\R)\) act diagonally through its usual action on projective lines, and define the orbit space with the quotient topology
\[
B_n=U_n/\operatorname{SL}_2(\R).
\]
The central kernel \(\{\pm I\}\) acts trivially, so this coarse orbit space is also the quotient by \(\operatorname{PSL}_2(\R)\). It is not silently replaced by a \(\operatorname{PGL}_2(\R)\) quotient, which would additionally allow orientation-reversing projective transformations.

The source asks, as a function of \(n\), to compute the cohomology of \(B_n\) and the number of its connected components. A standard precise cohomology target would be ordinary singular \(H^q(B_n;\Z)\) for all \(q\geq0\), with cup-product structure if intended. The source does not specify coefficients, support conditions, or whether a different quotient-theoretic cohomology is intended. Stack or equivariant cohomology is not the same as cohomology of the coarse orbit space, even though the central kernel does not change that orbit space. No Hausdorff or proper-action hypothesis is supplied merely by calling it a quotient. For \(n<3\), the image-size condition leaves no configurations, if those indices are included at all.

The supplied literature assessment checked 2026-08-21 says that the report defines this baby Teichm\"uller space but supplies no theorem resolving the cohomology and component questions; its open status is a dated triage assessment.
\subsection{Short English statement}
Compute the cohomology and connected-component count of the labeled cyclic configurations in \(\R\mathrm P^1\), with only adjacent collisions forbidden and at least three distinct values, modulo \(\operatorname{SL}_2(\R)\). Quotient and coefficient conventions matter.
\subsection{Sources}
\begin{itemize}
\item Local UnsolvedMath registry, permanent ID 30000510, catalogue entry, \emph{Topology of Baby Teichm\"uller Spaces}.
\item \emph{Teichm\"uller Space (Classical and Quantum)} (2007), Oberwolfach Reports, pp. 1537--1614, DOI \url{https://doi.org/10.4171/owr/2006/26}.
\item Supplied statement-verification claim and literature triage checked 2026-08-21 cite the official report; no separate cohomology or component theorem is supplied.
\end{itemize}
\end{document}
